
\documentclass{article}
\usepackage{colortbl}
\usepackage{makecell}
\usepackage{multirow}
\usepackage{supertabular}

\begin{document}

\newcounter{utterance}

\centering \large Interaction Transcript for game `hot\_air\_balloon', experiment `air\_balloon\_survival\_en\_reasoning off\_hard', episode 3 with Qwen3{-}32B{-}t0.0.
\vspace{24pt}

{ \footnotesize  \setcounter{utterance}{1}
\setlength{\tabcolsep}{0pt}
\begin{supertabular}{c@{$\;$}|p{.15\linewidth}@{}p{.15\linewidth}p{.15\linewidth}p{.15\linewidth}p{.15\linewidth}p{.15\linewidth}}
    \# & $\;$A & \multicolumn{4}{c}{Game Master} & $\;\:$B\\
    \hline

    \theutterance \stepcounter{utterance}  
    & & \multicolumn{4}{p{0.6\linewidth}}{
        \cellcolor[rgb]{0.9,0.9,0.9}{
            \makecell[{{p{\linewidth}}}]{
                \texttt{\tiny{[P1$\langle$GM]}}
                \texttt{You are participating in a collaborative negotiation game.} \\
\\ 
\texttt{Together with another participant, you must agree on a single set of items that will be kept. Each of you has your own view of how much each item matters to you (importance). You do not know how the other participant values the items. Additionally, you are given the effort each item demands.} \\
\texttt{You may only agree on a set if the total effort of the selected items does not exceed a shared limit:} \\
\\ 
\texttt{LIMIT = 3629} \\
\\ 
\texttt{Here are the individual item effort values:} \\
\\ 
\texttt{Item effort = \{"A21": 5, "B93": 793, "A11": 21, "B33": 842, "C08": 313, "B39": 842, "B96": 631, "C85": 226, "A82": 87, "B31": 761, "B20": 230, "A60": 287, "B09": 698, "C81": 641, "B21": 882\}} \\
\\ 
\texttt{Here is your personal view on the importance of each item:} \\
\\ 
\texttt{Item importance values = \{"A21": 138, "B93": 583, "A11": 868, "B33": 822, "C08": 783, "B39": 65, "B96": 262, "C85": 121, "A82": 508, "B31": 780, "B20": 461, "A60": 484, "B09": 668, "C81": 389, "B21": 808\}} \\
\\ 
\texttt{Goal:} \\
\\ 
\texttt{Your goal is to negotiate a shared set of items that benefits you as much as possible (i.e., maximizes total importance to YOU), while staying within the LIMIT. You are not required to make a PROPOSAL in every message {-} you can simply negotiate as well. All tactics are allowed!} \\
\\ 
\texttt{Interaction Protocol:} \\
\\ 
\texttt{You may only use the following structured formats in a message:} \\
\\ 
\texttt{PROPOSAL: \{'A', 'B', 'C', …\}} \\
\texttt{Propose keeping exactly those items.} \\
\texttt{REFUSE: \{'A', 'B', 'C', …\}} \\
\texttt{Explicitly reject opponent's proposal.} \\
\texttt{ARGUMENT: \{'...'\}} \\
\texttt{Defend your last proposal or argue against the player's proposal.} \\
\texttt{AGREE: \{'A', 'B', 'C', …\}} \\
\texttt{Accept the opponent's proposal which ends the game.} \\
\\ 
\\ 
\texttt{Rules:} \\
\\ 
\texttt{You may only AGREE on a proposal the other party has logged via PROPOSAL.} \\
\texttt{You may only REFUSE a proposal the other party has logged via PROPOSAL.} \\
\texttt{Total effort of any PROPOSAL or AGREE set must be ≤ LIMIT.} \\
\texttt{Do NOT reveal your hidden importance scores.} \\
\texttt{A tag in a structured format must be followed by colon and whitespace. The argument must be a python set containing 0 or more strings.} \\
\texttt{So, it must be of the form TAG: \{...\}} \\
\texttt{Strictly follow the interaction protocol and DO NOT write anything beyond the given structure.} \\
\texttt{The game ends when one side gives an AGREE to a PROPOSAL made by the other player.} \\
\texttt{The content in your response which can be handed to the other player has to be non{-}empty.} \\
\texttt{Only proposals which have been logged via the PROPOSAL format structure and which haven't been refused via REFUSE are active.} \\
\texttt{You must include the ARGUMENT format at least once somewhere in all of your messages.} \\
            }
        }
    }
    & \\ \\

    \theutterance \stepcounter{utterance}  
    & \multicolumn{4}{p{0.6\linewidth}}{
        \cellcolor[rgb]{0.9,0.9,0.9}{
            \makecell[{{p{\linewidth}}}]{
                \texttt{\tiny{[P1$\rangle$GM]}}
                \texttt{<think>} \\
\texttt{Okay, let's see. I need to negotiate a set of items that maximizes my total importance without exceeding the effort limit of 3629. First, I should calculate the total effort of all items to see if the sum is over the limit. Let me add them up.} \\
\\ 
\texttt{Looking at the item efforts: A21 is 5, B93 is 793, A11 is 21, B33 is 842, C08 is 313, B39 is 842, B96 is 631, C85 is 226, A82 is 87, B31 is 761, B20 is 230, A60 is 287, B09 is 698, C81 is 641, B21 is 882. Let me sum these up step by step.} \\
\\ 
\texttt{Starting with 0:} \\
\texttt{+5 = 5} \\
\texttt{+793 = 798} \\
\texttt{+21 = 819} \\
\texttt{+842 = 1661} \\
\texttt{+313 = 1974} \\
\texttt{+842 = 2816} \\
\texttt{+631 = 3447} \\
\texttt{+226 = 3673} \\
\texttt{+87 = 3760} \\
\texttt{+761 = 4521} \\
\texttt{+230 = 4751} \\
\texttt{+287 = 5038} \\
\texttt{+698 = 5736} \\
\texttt{+641 = 6377} \\
\texttt{+882 = 7259. } \\
\\ 
\texttt{Wow, that's way over the limit of 3629. So we can't take all items. Need to find a subset that's under 3629 and maximizes my importance.} \\
\\ 
\texttt{My importance values are: A21(138), B93(583), A11(868), B33(822), C08(783), B39(65), B96(262), C85(121), A82(508), B31(780), B20(4} \\
            }
        }
    }
    & & \\ \\

    \theutterance \stepcounter{utterance}  
    & & \multicolumn{4}{p{0.6\linewidth}}{
        \cellcolor[rgb]{0.9,0.9,0.9}{
            \makecell[{{p{\linewidth}}}]{
                \texttt{\tiny{[P1$\langle$GM]}}
                \texttt{Your response did not contain an argument. You must include the structured format ARGUMENT: \{...\} somewhere in your response. Try again} \\
            }
        }
    }
    & \\ \\

    \theutterance \stepcounter{utterance}  
    & \multicolumn{4}{p{0.6\linewidth}}{
        \cellcolor[rgb]{0.9,0.9,0.9}{
            \makecell[{{p{\linewidth}}}]{
                \texttt{\tiny{[P1$\rangle$GM]}}
                \texttt{<think>} \\
\texttt{Okay, let me try to figure this out. The user is in a negotiation game where they need to agree on a set of items with another participant. The goal is to maximize their own total importance without exceeding the effort limit of 3629. They provided their item efforts and importance values. } \\
\\ 
\texttt{First, I need to calculate the total effort of all items to see if it's over the limit. Adding up all the efforts: A21(5) + B93(793) + A11(21) + B33(842) + C08(313) + B39(842) + B96(631) + C85(226) + A82(87) + B31(761) + B20(230) + A60(287) + B09(698) + C81(641) + B21(882). Let me add these step by step. } \\
\\ 
\texttt{Starting with 5 + 793 = 798, +21 = 819, +842 = 1661, +313 = 1974, +842 = 2816, +631 = 3447, +226 = 3673, which is already over 3629. So even without adding the rest, the total is way over. So they can't take all items. } \\
\\ 
\texttt{Next, the user needs to select a subset. Since their importance values are given, they should prioritize items with high importance per unit effort. Let's calculate the ratio of importance to effort for each item. } \\
\\ 
\texttt{For example, A11 has importance 868 and effort 21, so ratio \textasciitilde{}41.33. B33 has 822/842 ≈ 0.976. C08 is 783/313 ≈ 2.5. B93 is 583/793 ≈ 0.735. A11 is a good candidate because of high ratio. Similarly, A82 has 508/87 ≈ 5} \\
            }
        }
    }
    & & \\ \\

    \theutterance \stepcounter{utterance}  
    & & \multicolumn{4}{p{0.6\linewidth}}{
        \cellcolor[rgb]{0.9,0.9,0.9}{
            \makecell[{{p{\linewidth}}}]{
                \texttt{\tiny{[P1$\langle$GM]}}
                \texttt{Your response did not contain an argument. You must include the structured format ARGUMENT: \{...\} somewhere in your response. Try again} \\
            }
        }
    }
    & \\ \\

    \theutterance \stepcounter{utterance}  
    & \multicolumn{4}{p{0.6\linewidth}}{
        \cellcolor[rgb]{0.9,0.9,0.9}{
            \makecell[{{p{\linewidth}}}]{
                \texttt{\tiny{[P1$\rangle$GM]}}
                \texttt{<think>} \\
\texttt{Okay, let me try to figure this out. The user is in a negotiation game where they need to agree on a set of items with another participant. The goal is to maximize their own total importance without exceeding the effort limit of 3629. They provided their item efforts and importance values. } \\
\\ 
\texttt{First, I need to calculate the total effort of all items to see if it's over the limit. Adding up all the efforts: A21(5) + B93(793) + A11(21) + B33(842) + C08(313) + B39(842) + B96(631) + C85(226) + A82(87) + B31(761) + B20(230) + A60(287) + B09(698) + C81(641) + B21(882). Let me add these step by step. } \\
\\ 
\texttt{Starting with 5 + 793 = 798, +21 = 819, +842 = 1661, +313 = 1974, +842 = 2816, +631 = 3447, +226 = 3673, which is already over 3629. So even without adding the rest, the total is way over. So they can't take all items. } \\
\\ 
\texttt{Next, the user needs to select a subset. Since their importance values are given, they should prioritize items with high importance per unit effort. Let's calculate the ratio of importance to effort for each item. } \\
\\ 
\texttt{For example, A11 has importance 868 and effort 21, so ratio \textasciitilde{}41.33. B33 has 822/842 ≈ 0.976. C08 is 783/313 ≈ 2.5. B93 is 583/793 ≈ 0.735. A11 is a good candidate because of high ratio. Similarly, A82 has 508/87 ≈ 5} \\
            }
        }
    }
    & & \\ \\

    \theutterance \stepcounter{utterance}  
    & & & \multicolumn{2}{p{0.3\linewidth}}{
        \cellcolor[rgb]{0.9,0.9,0.9}{
            \makecell[{{p{\linewidth}}}]{
                \texttt{\tiny{[GM$|$GM]}}
                \texttt{abort} \\
            }
        }
    }
    & & \\ \\

    \theutterance \stepcounter{utterance}  
    & & & \multicolumn{2}{p{0.3\linewidth}}{
        \cellcolor[rgb]{0.9,0.9,0.9}{
            \makecell[{{p{\linewidth}}}]{
                \texttt{\tiny{[GM$|$GM]}}
                \texttt{end game} \\
            }
        }
    }
    & & \\ \\

\end{supertabular}
}

\end{document}
